\bookchapter{Глава 44: Практика от 0 до 1 — полный цикл проекта}

\chapterintro{Теория пройдена — теперь пора за дело. Эта глава — не простое «повторяй за мной», а практическая методология «как думать, как спрашивать, как проверять». На трёх проектах разной сложности мы разберём полный ход размышлений — от требований до готового продукта. Каким бы инструментом вы ни пользовались, этот процесс универсален.}

\booksection{Проект первый: минималистичный сайт личных инструментов (начальный уровень)}

\booksubsection{Обдумываем требования}

Цель: создать веб-страницу — личный набор инструментов, показывающий входы к 3–4 полезным мини-инструментам. Это чисто фронтенд-проект: бэкенд и база данных не нужны. Главное здесь не сложность самих инструментов, а опыт прохождения полного пути «от требований до готового продукта».

Пять элементов требований: Кто (Who) — для личного пользования. Что (What) — веб-страница, показывающая входы к инструментам. Зачем (Why) — собрать часто используемые инструменты на одной странице для удобного доступа. Как (How) — один файл HTML+CSS+JS, без фреймворков. Ограничение (Constraint) — уложиться в 30 минут, без развёртывания.

\booksubsection{Отдаём инструкцию}

Согласно принципу Think Before Coding («сначала подумай, потом кодируй»), ваша инструкция должна содержать достаточно информации, чтобы ИИ понял цель, но не обязана расписывать каждую деталь реализации:

\begin{lstlisting}
Создайте веб-страницу личного набора инструментов, требования:
1. HTML одним файлом (CSS и JS инлайном), без фреймворков
2. Тёмная тема, современный минимализм
3. Четыре карточки-инструмента:
   - Суммаризатор текста
   - Счётчик слов
   - Пипетка цветов
   - Генератор QR-кодов
4. У каждой карточки иконка, заголовок и короткое описание
5. Анимация при наведении на карточки
6. Адаптивная вёрстка, удобна на телефоне
7. Сверху полоса с заголовком
Сразу дайте полный HTML-файл.
\end{lstlisting}

\booksubsection{Приёмка и итерации}

Когда ИИ выдаст код, проведите приёмку по трёхуровневой проверке. Первый уровень — проверка функциональности: открывается ли файл в браузере? Отображаются ли все 4 карточки? Работает ли анимация при наведении? Второй уровень — проверка границ: нормально ли выглядит макет при размере экрана телефона? Переносятся ли карточки на новую строку при узком окне? Третий уровень — проверка архитектуры: читаем ли код? Собран ли CSS внутри одного тега style? Нет ли явной избыточности?

Если захотите реализовать и функциональность какой-то карточки (а не только вход к ней), добавляйте её постепенно — за один раз добавляйте функциональность одного инструмента, проверяйте и лишь затем переходите к следующему. Это практическое применение стратегии «мелкие шаги и быстрый ритм».

\booksection{Проект второй: fullstack-система учёта финансов (средний уровень)}

\booksubsection{Обдумывание требований и PRD}

Цель: fullstack-приложение для учёта финансов с фронтендом, бэкендом и базой данных. Этот проект заметно сложнее первого — здесь пригодится методология работы с требованиями из главы 40.

Сначала напишите PRD. По шаблону из третьего раздела главы 40 опишите обзор проекта, список функций, техническое решение, модель данных и план разработки. Выбор технологий должен опираться на принцип Simplicity First («простота прежде всего») — на фронтенде берём React, но без Redux и других сложных средств управления состоянием; на бэкенде — Node.js, но без микросервисов; база данных — SQLite, а не PostgreSQL (на этапе разработки).

План разработки разбит на 5 шагов: Шаг 1 — каркас API на бэкенде (сервер Express + SQLite + CRUD-эндпоинты); Шаг 2 — базовые страницы фронтенда (список записей + форма добавления); Шаг 3 — связка фронтенда с бэкендом (вызовы API + показ данных); Шаг 4 — категории и фильтрация; Шаг 5 — месячная сводка и простые диаграммы. После каждого шага на руках «работающая вещь».

\booksubsection{Настраиваем файл памяти проекта}

Прежде чем начать кодировать, создайте файл памяти проекта. Это ключевая практика контекстной инженерии — ИИ в начале каждой беседы будет знать предысторию проекта.

\begin{lstlisting}
# Проект: полностековое приложение финансов
## Обзор
Приложение личных финансов: записи доходов/расходов, категории, месячные итоги.
## Технологии
- Фронтенд: React + TypeScript + Vite
- Бэкенд: Node.js + Express + TypeScript
- База данных: SQLite (через better-sqlite3)
- Общие определения типов для фронта и бэка
## Соглашения API
- Базовый путь: /api
- Формат ответа:{ success: boolean, data: any, error: string }
- Стиль RESTful
## Модель данных
Record { id, amount, type, category, note, date }
## Стиль кода
- Строгий режим TypeScript
- Отступ в 2 пробела
- У функций есть JSDoc-комментарии
- Комментарии и коммиты на русском
\end{lstlisting}

\booksubsection{Реализуем шаг за шагом}

Следуя плану разработки, выполняйте каждый шаг в отдельной беседе. Вот пример инструкции для первого шага:

\begin{lstlisting}
Нужен каркас серверного API.
Требования:
1. Сервер Express, порт 3001
2. База данных SQLite, таблица records
3. Поля таблицы records:
   id (TEXT PRIMARY KEY)
   amount (REAL)
   type (TEXT: 'income'|'expense')
   category (TEXT)
   note (TEXT)
   date (TEXT: YYYY-MM-DD)
4. CRUD-эндпоинты:
   GET /api/records — получить все записи
   POST /api/records — добавить запись
   PUT /api/records/:id — обновить запись
   DELETE /api/records/:id — удалить запись
5. Мидлварь обработки ошибок
Дайте полную структуру серверного кода.
\end{lstlisting}

Обратите внимание на особенности этой инструкции: она описывает только текущий шаг (каркас бэкенда) и не подмешивает требования фронтенда; она явно перечисляет поля и эндпоинты (задаёт ИИ чёткие рамки), оставляя конкретную реализацию на его усмотрение; она следует принципу Goal-Driven Execution — говорит «что делать», а не «как делать».

Как только ИИ выдаст код, сразу проверяйте: запустите сервер и протестируйте каждый эндпоинт через curl или Postman. Убедившись, что всё работает, открывайте новую беседу для шага 2 (базовые страницы фронтенда). Каждый шаг проходит цикл «инструкция — проверка — новая беседа».

\booksection{Проект третий: приложение с усилением ИИ (продвинутый уровень)}

\booksubsection{Замысел проекта}

Цель: приложение, подключающееся к API большой языковой модели, — например «ИИ-помощник по учёбе»: пользователь вводит учебный материал, а ИИ делает конспект, связывает понятия и советует учебные ресурсы. Особенность проекта в том, что ваше приложение само вызывает ИИ, а значит, вам предстоит решать новые задачи: вызовы API, промпт-инженерия, обработка ошибок.

Самый важный урок этого проекта: когда вы «с помощью ИИ пишете код, который вызывает ИИ», к промпт-инженерии нужно относиться особенно внимательно. Промпт, с которым ваше приложение обращается к API большой модели, тоже должен подчиняться принципам из этого руководства — ясность, лаконичность, ориентация на цель.

\booksubsection{Ключевые технические моменты}

В этом проекте появятся несколько технических моментов, которых не было в первых двух:

\begin{itemize}
  \item Управление API-ключами: ключи ни в коем случае нельзя зашивать в код — управляйте ими через переменные окружения;
  \item Обработка потоковых ответов: ответы больших моделей обычно приходят потоком, и текст нужно показывать по мере поступления;
  \item Контроль расхода токенов: большие модели тарифицируются по токенам, поэтому расход нужно ограничивать на уровне приложения;
  \item Управление шаблонами промптов: для разных функций нужны разные шаблоны промптов, и их стоит вести упорядоченно.
\end{itemize}

Все эти моменты можно поручить ИИ, но вы должны понимать «почему» за каждым: почему ключи нельзя зашивать в код? зачем нужны потоковые ответы? почему расход токенов нужно ограничивать? Понимая «почему», вы сумеете оценить, разумен ли предложенный ИИ план.

\booksection{Универсальный чек-лист рабочего процесса}

Независимо от масштаба проекта ниже — универсальный чек-лист практики программирования с ИИ. Берясь за новый проект, пройдитесь по нему, чтобы не упустить ключевые шаги.

\begin{center}
\begin{tabularx}{\linewidth}{|X|X|X|}
\hline
\textbf{Этап} & \textbf{Что делаем} & \textbf{Соответствующий принцип} \tabularnewline
\hline
1. Обдумывание требований & Пройти пять элементов требований и ясно понять, что делаем & Think Before Coding \tabularnewline
2. Разбиение требований & Разделить крупное требование на малые шаги, проверяемые по отдельности & Simplicity First \tabularnewline
3. Написать PRD & Оформить в структурированный документ на память & Goal-Driven \tabularnewline
4. Настроить файл памяти & Создать CLAUDE.md или похожий файл & Harness Engineering \tabularnewline
5. Реализация по шагам & Один шаг — одна беседа: инструкция — проверка — следующий шаг & Surgical Changes \tabularnewline
6. Тестирование и проверка & Сразу после каждого шага проверять правильность работы & Iterate \& Verify \tabularnewline
7. Ревью кода & Проверить разумность архитектуры, безопасность, сопровождаемость & Iterate \& Verify \tabularnewline
8. Итерации и улучшения & Дорабатывать по итогам тестов и ревью & Iterate \& Verify \tabularnewline
\hline
\end{tabularx}
\end{center}

\begin{center}\small\color{gray}Универсальный чек-лист практики программирования с ИИ\end{center}

\begin{tcbpractice}{Попробуйте сами: за эту неделю закончите «проект первый»}
Возьмите «сайт личных инструментов» из первого раздела главы 44 как своё первое домашнее задание по программированию с ИИ: потратьте 5 минут на пять элементов требований, отправьте инструкцию ИИ (подойдёт любой инструмент) и примите работу по трёхуровневой проверке. Завершите за 30 минут — и ваша первая работа в программировании с ИИ готова. Обязательно сохраните её: это памятка о том, как вы превратились из «читателя» в «практика».

\end{tcbpractice}

\begin{tcbsummary}{Краткие итоги главы}
\begin{itemize}
  \item Проект 1 (начальный): чисто фронтендовый сайт инструментов — полный цикл, делается за 30 минут
  \item Проект 2 (средний): fullstack-система учёта финансов — PRD + файл памяти + 5 шагов поочерёдно, после каждого всё «работает»
  \item Проект 3 (продвинутый): ИИ-помощник по учёбе — ИИ пишет код, вызывающий ИИ; внимание к промпт-инженерии
  \item Универсальный процесс из 8 шагов: обдумывание требований разбиение PRD файл памяти реализация по шагам тесты ревью итерации
  \item Каждый шаг проходит цикл «инструкция — проверка — новая беседа»
\end{itemize}

\end{tcbsummary}
